\section{Conclusion}
\label{sec:conclusion}

We introduced CLP-CSGM Ridge, a conditional generative compressed sensing method
for petrophysical property estimation from dense logs and sparse core
measurements. The method learns a target-property decoder, predicts a
log-conditioned latent prior with ridge regression, and refines the latent code
at test time by enforcing sparse measurement consistency.

In cross-well low-data Vc prediction, CLP-CSGM Ridge outperformed the strongest
direct \([u,b]\) autoencoder baseline in all tested training-size and
measurement-ratio cells. In the real-well F03 GR-only porosity case, it
achieved the best average RMSE, with its clearest advantage in the sparsest
measurement regimes. These results indicate that conditional latent generative
priors are a promising alternative to fixed-basis sparse residual recovery for
petrophysical sparse-measurement assimilation.

Future work should evaluate the method on additional fields, compare richer
conditional priors, investigate uncertainty quantification in latent space, and
study measurement designs beyond coordinate subsampling.
